\section{Experimental Setup}
\label{sec:setup}

\subsection{Dataset and Model}

We evaluate on the RadioML 2016.10a dataset~\cite{otoole2016rml}, a widely
used benchmark for automatic modulation classification comprising 220{,}000 IQ
signal frames spanning 11 modulation classes: eight digital (BPSK, QPSK, 8PSK,
QAM16, QAM64, PAM4, CPFSK, GFSK) and three analog (AM-DSB, AM-SSB, WBFM).
Each frame consists of 128 complex-valued samples represented as a $2\times128$
real tensor of in-phase (I) and quadrature (Q) components.  The dataset covers
SNR values from $-20$\,dB to $+18$\,dB in 2\,dB increments, yielding 20 SNR
levels and 1{,}000 frames per (modulation, SNR) cell.

We partition the data using stratified sampling that preserves the joint
(modulation, SNR) distribution across splits: 70\% training, 15\%
validation, and 15\% test.  All reported accuracies are computed on the
held-out test split.

The classifier under attack is the Adaptive Wavelet Network (AWN)~\cite{li2023awn},
a pretrained model that achieves 92.6\% classification accuracy at SNR $\geq 0$\,dB
on the clean test set.  We use the publicly available pretrained checkpoint without
any adversarial fine-tuning, so the model represents a realistic deployment target
that has not been hardened against adversarial inputs.

\subsection{Attack Configurations}

We evaluate five attacks drawn from two families: single-step gradient attacks
(FGSM), iterative gradient attacks (PGD), and optimization-based attacks
(CW-L2, EAD-L1, EAD-EN). All attacks are implemented via the torchattacks
library~\cite{kim2020torchattacks} using a \texttt{Model01Wrapper} that maps
IQ signals into the $[0,1]$ range expected by torchattacks and returns logits
only.  We use per-sample min-max normalization (\texttt{ta\_box=minmax}) to
map each sample to $[0,1]$ so that epsilon is expressed relative to the
sample's own dynamic range---a natural choice for IQ signals whose peak-to-peak
amplitude varies widely across modulations and SNR levels.

\begin{table}[t]
\caption{Attack configurations. All attacks target untargeted misclassification
  in a white-box setting. Epsilon values are expressed after per-sample min-max
  normalization to $[0,1]$.}
\label{tab:attacks}
\centering
\small
\begin{tabular}{llll}
\toprule
Attack & Norm & Key Parameters & Source \\
\midrule
FGSM   & $\ell_\infty$ & $\varepsilon = 0.03$                     & \cite{goodfellow2015fgsm} \\
PGD    & $\ell_\infty$ & $\varepsilon = 0.03$, steps $= 10$, $\alpha = 0.01$ & \cite{madry2018pgd} \\
CW     & $\ell_2$      & $c = 1.0$, steps $= 100$, lr $= 0.01$   & \cite{carlini2017towards} \\
EAD-L1 & $\ell_1$      & $c_0 = 0.01$, $\beta = 0.01$, steps $= 100$ & \cite{chen2020ead} \\
EAD-EN & elastic net   & $c_0 = 0.01$, $\beta = 0.01$, steps $= 100$ & \cite{chen2020ead} \\
\bottomrule
\end{tabular}
\end{table}

Table~\ref{tab:attacks} summarizes the attack hyperparameters.  For FGSM and
PGD, $\varepsilon = 0.03$ after min-max normalization corresponds to roughly
60\% of the average per-sample dynamic range at high SNR---a moderate threat
that leaves the signal recognizable but is effective enough to reduce accuracy
substantially.  For CW and EAD, the optimization constant $c$ controls the
trade-off between perturbation minimization and attack success; we use $c=1.0$
and $c_0=0.01$, respectively, as these yield near-zero undefended accuracy on
most modulations at high SNR while keeping perturbation norms in a physically
plausible range.

\subsection{Defense Parameters}

We evaluate nine defense configurations: two proposed adaptive spectral methods
(Adaptive-K, Spectral-Gated) and seven baselines (Kalman, Wiener,
Savitzky-Golay, Gaussian, FIR, and Randomized Smoothing), plus the undefended
baseline.  The proposed methods operate entirely in the frequency domain without
retraining.

\begin{table}[t]
\caption{Defense configurations. Classical filter parameters are calibrated
  per-SNR via grid search on the validation set; representative values at
  SNR $= 10$\,dB are shown. Proposed methods (Adaptive-K and Spectral-Gated)
  require no per-SNR calibration.}
\label{tab:defenses}
\centering
\small
\begin{tabular}{llp{4.5cm}}
\toprule
Defense & Type & Key Parameters (SNR $= 10$\,dB) \\
\midrule
Adaptive-K      & Proposed & $K$ auto per sample; energy threshold $\eta=0.95$ \\
Spectral-Gated  & Proposed & $K=50$; spectral flatness threshold $= 0.5$ \\
Kalman          & Classical & $Q=0.005$, $R=0.001$; CPU sequential \\
Wiener          & Classical & noise $= 0.01$, filter len $= 3$; freq.\ domain \\
Savitzky-Golay  & Classical & window $= 5$, order $= 1$; time domain \\
Gaussian        & Classical & $\sigma=0.5$; GPU depthwise conv1d \\
FIR             & Classical & cutoff $= 0.15$, taps $= 63$; GPU depthwise conv1d \\
Rand.\ Smoothing & Baseline  & $\sigma=0.1$, $k=1$ copies; certified robustness \\
\bottomrule
\end{tabular}
\end{table}

Table~\ref{tab:defenses} summarizes the defense parameters.  For the five
classical filters (Kalman, Wiener, Savitzky-Golay, Gaussian, FIR), we perform
a grid search over the parameter space at each SNR level using the validation
set with CW attack.  The best-performing parameter combination per SNR is stored
in a calibration file (\texttt{calibration\_params.json}) and loaded at test
time, giving the classical baselines every possible advantage.  Randomized
Smoothing~\cite{cohen2019certified} adds Gaussian noise ($\sigma=0.1$) before
classification using a single forward pass ($k=1$), consistent with its use
as an inference-time defense rather than a certified smoothed classifier.

\subsection{Evaluation Metrics}

We report the following metrics:

\textbf{Per-SNR accuracy}: fraction of test samples correctly classified at
each SNR level $s \in \{0, 2, \ldots, 18\}$ dB.  We focus on non-negative SNR
because at $\text{SNR} < 0$\,dB the clean accuracy is below 50\% and adversarial
perturbations are masked by channel noise.

\textbf{Weighted average accuracy}: the mean per-SNR accuracy, weighted by the
number of test samples at each SNR level.  Since RML2016.10a is balanced, this
reduces to the simple mean over the 10 SNR levels $\geq 0$\,dB evaluated.

\textbf{Perturbation budget curves}: accuracy as a function of attack strength
($\varepsilon$ for FGSM/PGD; optimization constant $c$ for CW/EAD) to assess
defense robustness across the threat spectrum.

\textbf{Confusion matrices}: per-class classification accuracy before and after
defense at a fixed SNR (18\,dB) to identify which modulation classes benefit
most from the proposed methods.

All evaluations use 200 samples per (SNR, modulation) cell from the test split,
yielding 2{,}200 samples per SNR level (200 samples $\times$ 11 modulation classes).
